\begin{tablalafrox}{Perfiles operativos y modalidad de capacitación. Fuente: (Subdocumento 3, Anexo 3.I) y (Bases Técnicas del caso, secciones 2.4 y 8, pp. 5, 13--17).}{tab:T18-perfiles}
  {L{3.5cm} L{3.4cm} Y}
  {\cab{Perfil} & \cab{Personas} & \cab{Cómo se capacita sin detener la operación}}
  Preparadores Talca & 120 & Aprendizaje en el puesto de 2 horas como máximo (RNF-05.03) y un tutor por turno; capacitación continua por la rotación del 38 \% \\
  Preparadores Concepción & Hasta 60, según el supuesto S-39 & Aprendizaje en el puesto de hasta 2 horas (RNF-05.03), tutor por turno y certificación antes del cierre de la ola. Confirmar nómina, turnos y cobertura con el levantamiento del mes 1. \\
  Preventistas & 62 & Prueba de captura sin señal en su propia ruta y acompañamiento en ruta \\
  Tripulación propia y peonetas & 84 & Prueba de entrega digital y rendición antes de la ola, y acompañamiento en ruta \\
  Conductores de transportistas & ≈ 160 & Incorporación en el andén al asignarse a una ruta, con autenticación y vinculación persona–vehículo–viaje (S-05) \\
\end{tablalafrox}
